\documentclass[11pt,a4paper]{article}
\usepackage[utf8]{inputenc}
\usepackage[T1]{fontenc}
\usepackage{amsmath,amssymb,amsthm,mathtools}
\usepackage{geometry}
\geometry{margin=1in}
\usepackage{graphicx}
\usepackage{float}
\usepackage{hyperref}
\usepackage{cleveref}
\usepackage{booktabs}
\usepackage{enumitem}
\usepackage{caption}
\usepackage{subcaption}
\usepackage{algorithm}
\usepackage{algpseudocode}
\usepackage[numbers]{natbib}   % [numbers] added: plain style -> numeric citations
\bibliographystyle{plain}
\newtheorem{theorem}{Theorem}[section]
\newtheorem{lemma}[theorem]{Lemma}
\newtheorem{proposition}[theorem]{Proposition}
\newtheorem{corollary}[theorem]{Corollary}
\theoremstyle{definition}
\newtheorem{definition}[theorem]{Definition}
\newtheorem{remark}[theorem]{Remark}
\newcommand{\R}{\mathbb{R}}
\newcommand{\Z}{\mathbb{Z}}
\newcommand{\N}{\mathbb{N}}
\newcommand{\E}{\mathbb{E}}
\newcommand{\Var}{\mathrm{Var}}
\newcommand{\Cov}{\mathrm{Cov}}
\newcommand{\tr}{\mathrm{tr}}
\newcommand{\dd}{\,\mathrm{d}}
\newcommand{\avg}[1]{\langle #1 \rangle}
% --- added: macros used in the body but absent from the provided preamble ---
\numberwithin{equation}{section}   % added: per-section equation numbers (matches theorem[section])
\newcommand{\bI}{\mathbf{I}}
\newcommand{\bphi}{\boldsymbol{\varphi}}
\newcommand{\Lgen}{\mathcal{L}}
\newcommand{\Ladj}{\mathcal{L}^{*}}

\title{Numerical study of an energy cascade model derived from a dispersive equation}
\author{Jayleen Jiang \and Yao Li\thanks{Department of Mathematics and Statistics, University of Massachusetts Amherst, Amherst, MA 01003, USA. \texttt{liyao@umass.edu}}}
\date{\today}

\begin{document}
\maketitle


\section{Introduction}

The transfer of energy between the modes of a nonlinear Hamiltonian system is a
basic question across the sciences, from fluids and oceanography to climate.
When external forcing and dissipation act on a few modes while the bulk---the
\emph{inertial range}---evolves intrinsically, the natural framework is
out-of-equilibrium statistical physics; for wave systems it is \emph{wave} (or
\emph{weak}) \emph{turbulence theory} \cite{Nazarenko}, in which energy is
transported between Fourier modes, that is, across spatial scales. The classical
description proceeds through the wave kinetic equation, whose power-law
stationary solutions are the Kolmogorov--Zakharov cascade spectra
$n(k)=|k|^{-d}$ and $n(k)=|k|^{-d+2/3}$ (the direct and inverse cascades). These
spectra are themselves nonequilibrium steady states (NESS), in which energy
flows at constant flux through the inertial range; their rigorous justification
remains largely open \cite{HLNS}.

Following \cite{HLNS}, we adopt a complementary viewpoint: rather than the
kinetic equation, we study the invariant measure---the NESS itself---of a
forced and dissipated reduction of the nonlinear Schr\"odinger (NLS) equation.
The present paper is the numerical companion to that program. We compute and
analyze the NESS of the energy cascade chain: for the three-mode system treated
analytically in \cite{HLNS}, and, for the first time, for long chains, where the
analysis is open.

\subsection{Derivation of the model}\label{sec:derivation}

Consider the cubic NLS on a torus,
\begin{equation}\label{eq:nls}
  -i\partial_t v + \Delta v = \pm |v|^2 v .
\end{equation}
In Fourier variables $a_K$ the cubic nonlinearity couples wavevectors with
$K_1-K_2+K_3=K$. Retaining only the \emph{resonant} interactions, those with
$|K_1|^2-|K_2|^2+|K_3|^2=|K|^2$, gives the resonant NLS system, which
approximates \eqref{eq:nls} for sufficiently small data. Colliander, Keel,
Staffilani, Takaoka and Tao \cite{CKSTT} restrict this system to a finite set of
frequencies
\[
  \Lambda=\Lambda_1\cup\Lambda_2\cup\cdots\cup\Lambda_n\subset\mathbb{Z}^2,
\]
on which the amplitude is constant within each generation $\Lambda_j$. The
generations are produced by an explicit combinatorial ``nuclear family''
construction that embeds successive $\Lambda_j$ in $\mathbb{Z}^2$ through exact
rational (Pythagorean) rotations, with $\Lambda_1$ supported at low and
$\Lambda_n$ at high frequency. Writing $c_j(t)$ for the common amplitude on
$\Lambda_j$, the resonant system collapses to the \emph{toy model} of
\cite{CKSTT},
\begin{equation}\label{eq:toy}
  -i\dot c_j = 2\Big(\sum_k |c_k|^2\Big)c_j - |c_j|^2 c_j
  + 2\big(c_{j-1}^2 + c_{j+1}^2\big)\bar c_j,
  \qquad j=1,\dots,n,\quad c_0=c_{n+1}=0 .
\end{equation}
In \cite{CKSTT} this model is shown to admit an energy-cascade solution whose
mass, concentrated on $\Lambda_1$ at $t=0$, migrates to $\Lambda_n$ at a later
time---a deterministic route to the growth of Sobolev norms of \eqref{eq:nls}.

System \eqref{eq:toy} is Hamiltonian. In action--angle coordinates
$c_j=\sqrt{I_j}\,e^{i\varphi_j}$, with the \emph{bond angles}
$\theta_k:=2(\varphi_k-\varphi_{k-1})$, it is generated by the Hamiltonian
\begin{equation}\label{eq:hamiltonian}
  H(\bI,\bphi) \;=\; \Big(\sum_{k=1}^{n} I_k\Big)^{2}
  \;-\;\tfrac12\sum_{k=1}^{n} I_k^{2}
  \;+\; 2\sum_{k=2}^{n} I_{k-1} I_k \cos\theta_k ,
\end{equation}
with the Hamiltonian equations of motion
\begin{equation}\label{eq:ham-flow}
  -\partial_{\varphi_j}H = 4 I_j\big(I_{j-1}\sin\theta_j - I_{j+1}\sin\theta_{j+1}\big),
  \qquad
  \partial_{I_j}H = 2M - I_j + 2 I_{j-1}\cos\theta_j + 2 I_{j+1}\cos\theta_{j+1},
\end{equation}
where $M=\sum_k I_k$ and $I_0=I_{n+1}=0$.

\begin{remark}[Normalization convention]\label{rem:convention}
The Hamiltonian \eqref{eq:hamiltonian} is twice the action--angle Hamiltonian of
\cite{HLNS}; equivalently the drift \eqref{eq:ham-flow} is twice theirs, and the
equilibrium measure below carries the factor $1/2T$ rather than $1/T$. This
amounts to a rescaling of time and is fixed throughout so that every formula
matches the simulation code that produces our numerical data. For $n=3$ with the
middle mode as reference one writes $\theta_1=2(\varphi_1-\varphi_2)$,
$\theta_3=2(\varphi_3-\varphi_2)$, recovering the reduced coordinates of
\cite{HLNS}.
\end{remark}

To realize an energy cascade as a statistical steady state one injects energy at
the low-frequency end and extracts it at the high-frequency end: the two end
modes $j=1$ and $j=n$ are coupled to stochastic heat baths at temperatures $T_1$
and $T_n$. This produces the It\^o system
\begin{align}\label{eq:sde-bulk}
  \dd I_j   &= -\,\partial_{\varphi_j} H \dd t
             \;+\; b^{I}_j \dd t \;+\; \sigma^{I}_j \dd W^{I}_j, \\
  \dd \varphi_j &= \phantom{-}\partial_{I_j} H \dd t
             \;+\; b^{\varphi}_j \dd t \;+\; \sigma^{\varphi}_j \dd W^{\varphi}_j,
             \qquad j = 1,\dots,n, \nonumber
\end{align}
in which the bath drift and diffusion $(b,\sigma)$ are supported only on
$j\in\{1,n\}$. We take the fluctuation--dissipation (Langevin) form: at the hot
end $j=1$,
\begin{equation}\label{eq:bath}
  b^{I}_1 = 2\gamma\big(2T_1 - I_1\,\partial_{I_1}H\big),\quad
  \sigma^{I}_1 = 2\sqrt{2\gamma T_1 I_1},\quad
  b^{\varphi}_1 = 2\gamma I_2 \sin\theta_2,\quad
  \sigma^{\varphi}_1 = \sqrt{2\gamma T_1/I_1},
\end{equation}
and symmetrically at the cold end $j=n$ with $T_1\to T_n$; here $\gamma>0$ is the
coupling strength. The bath \eqref{eq:bath} is constructed so that at equilibrium
$T_1=T_n=T$ the system admits the explicit Gibbs invariant density
\begin{equation}\label{eq:gibbs}
  \rho_{\mathrm{eq}}(\bI,\bphi) \;\propto\; \exp\!\big(-H(\bI,\bphi)/2T\big).
\end{equation}

\begin{remark}[Choice of thermostat]\label{rem:bath}
The analysis of \cite{HLNS} uses a different thermostat---dissipation of the form
$\gamma(T_j-I_j^3)$ with a tailored multiplicative angle noise---chosen to make
the Lyapunov construction tractable; under that choice the Gibbs measure
\eqref{eq:gibbs} is \emph{not} invariant, even at equal temperatures. For our
numerical study we instead use the fluctuation--dissipation bath
\eqref{eq:bath}, whose equilibrium measure is the explicit Gibbs density
\eqref{eq:gibbs}; this provides a stringent check on the Fokker--Planck solver
(\cref{sec:fp-numerics}). The reduced-system analysis of the stabilization
mechanism in \cref{sec:stabilization} retains the original $\gamma(T-I^3)$ bath
of \cite{HLNS} for direct comparison with the proof.
\end{remark}

\subsection{Known results and open problems}\label{sec:known}

For the three-mode chain ($n=3$), in which only the middle mode $I_2$ is
unforced, Hani, Li, Nahmod and Staffilani \cite{HLNS} establish the existence,
uniqueness and ergodicity of the NESS: when the two bath temperatures are
sufficiently different, the system possesses a unique invariant probability
measure $\pi$, and the law of the process converges to $\pi$ in total variation
at a power-law rate in time. The proof rests on a novel construction of a
Lyapunov function through an elliptic Feynman--Kac equation. The central
difficulty is that, unlike microscopic heat conduction where energy simply flows
from hot to cold, here the transfer of energy to and from the unforced middle
mode depends sensitively on the phase variables being ``correct'': the system
must be prevented both from \emph{overheating} (energy trapped in $I_2$) and from
\emph{freezing} (energy blocked at $I_2$), and which of these threatens depends on
whether $I_2$ is large or small. This phase-dependent stabilization mechanism is
the analytical heart of \cite{HLNS} and the object of our \cref{sec:stabilization}.

Two limitations of the analytical theory motivate the present work. First, the
results of \cite{HLNS} are confined to $n=3$; the authors explicitly leave longer
chains open, noting that already the three-mode case is demanding. Second, the
NESS is not given in closed form, so quantitative properties---energy profiles,
transport coefficients, local statistics, the relaxation spectrum---are
inaccessible to the existence theory. Both gaps are naturally addressed
numerically.

\subsection{Contributions and organization}\label{sec:contributions}

This paper is a numerical study of the NESS of the energy cascade model
\eqref{eq:sde-bulk}. After collecting the thermodynamic and spectral framework in
\cref{sec:thermo,sec:fp,sec:methods}, we present two groups of results.

\Cref{sec:long-chain} treats the \emph{long chain}, the regime left open by
\cite{HLNS}. We characterize the nonequilibrium steady state---the energy profile
and a local thermodynamic equilibrium of local-Gibbs type for adjacent modes---
and we study heat transport, finding an anomalous, sub-diffusive scaling of the
current with chain length (an insulating chain, opposite to the
Fermi--Pasta--Ulam behavior \cite{LepriLiviPoliti}), together with the
associated entropy production.

\Cref{sec:short-chain} returns to the \emph{three-mode chain} for a detailed
study. We solve the stationary Fokker--Planck equation by a mesh-free,
data-driven neural-network method \cite{ZhaiDobsonLi}, validating it against the
exact Gibbs measure \eqref{eq:gibbs} in equilibrium; we then examine the
nonequilibrium steady state, the phase-locking stabilization mechanism, and the
spectrum of the generator, which we find to be overdamped at the simulated
coupling. The numerical tools---a high-throughput SIMD Monte Carlo integrator, a
kernel density estimator for the reference measure, and a trajectory-based
backward solver for the dominant eigenpair---are described as they arise.


\section{Preliminary}

This section collects the thermodynamic observables we measure for the long
chain (\cref{sec:thermo}), the spectral (Fokker--Planck) framework for the
short chain (\cref{sec:fp}), and the numerical methods used to solve the
stationary Fokker--Planck equation and to extract eigenfunctions of the
generator (\cref{sec:methods}).

\subsection{Nonequilibrium thermodynamics of the chain}\label{sec:thermo}

The Hamiltonian flow \eqref{eq:ham-flow} conserves the total mass
$M=\sum_k I_k$, and the evolution of each action obeys a discrete continuity
equation
\begin{equation}\label{eq:continuity}
  \dot I_j \;=\; J_j - J_{j+1},
  \qquad
  J_j \;:=\; 4\, I_{j-1} I_j \sin\theta_j ,
\end{equation}
so that $J_j$ is the \emph{action current} flowing from mode $j-1$ into mode
$j$ (with $J_1=J_{n+1}=0$ for the closed chain). In the nonequilibrium steady
state (NESS) the expected actions are stationary, $\E[\dot I_j]=0$, hence the
mean current is uniform along the chain, $\E[J_j]\equiv \E[J]$ for all interior
$j$.

\begin{definition}[Thermal conductivity and Fourier's law]\label{de:fourier}
Let $\nabla T$ denote the effective temperature gradient sustained across a
chain of length $n$ by baths at $T_1\neq T_n$. The chain obeys Fourier's law
with finite conductivity $\kappa$ if
\begin{equation}\label{eq:fourier}
  \E[J] \;=\; -\,\kappa\,\nabla T,
  \qquad
  \kappa \;=\; \lim_{n\to\infty}\kappa_n \in (0,\infty).
\end{equation}
Anomalous transport corresponds to $\kappa_n\sim n^{\alpha}$ with $\alpha\neq0$:
$\alpha>0$ is superdiffusive and $\alpha<0$ (so $\kappa_n\to0$) signals an
insulator.
\end{definition}

Two further observables characterize the NESS. The \emph{entropy production
rate} associated with the two heat baths is
\begin{equation}\label{eq:entropy-prod}
  \dot S \;=\; -\frac{\E[\dot Q_1]}{T_1} - \frac{\E[\dot Q_n]}{T_n} \;\ge\; 0,
\end{equation}
where $\dot Q_1,\dot Q_n$ are the mean heat fluxes delivered by the baths; in
the NESS the energy balance forces $\E[\dot Q_1]=-\E[\dot Q_n]=:\E[\dot Q]$ and
$\dot S = \E[\dot Q]\,(1/T_n - 1/T_1)$. The second law $\dot S\ge 0$ is refined
by the steady-state \emph{fluctuation theorem} of Gallavotti--Cohen
\cite{GallavottiCohen}: writing $\Sigma_\tau$ for the entropy produced over a
window of length $\tau$, its large-deviation rate function $e(\cdot)$ satisfies
\begin{equation}\label{eq:ft}
  \lim_{\tau\to\infty}\frac{1}{\tau}\,
  \ln\frac{\mathrm{Prob}(\Sigma_\tau/\tau = a)}{\mathrm{Prob}(\Sigma_\tau/\tau = -a)} \;=\; a .
\end{equation}
We test \eqref{eq:fourier}, \eqref{eq:entropy-prod} and \eqref{eq:ft}
numerically in \cref{sec:thermal-conductivity,sec:entropy}.

\subsection{Fokker--Planck equation and spectral decomposition}\label{sec:fp}

Write $x=(\bI,\bphi)$ and let the system \eqref{eq:sde-bulk} have drift $f(x)$
and diffusion matrix $\Sigma(x)=\sigma(x)\sigma(x)^{\!\top}$. The probability
density $u(x,t)$ evolves by the (forward) Fokker--Planck, or Kolmogorov forward,
equation $\partial_t u = \Ladj u$, where
\begin{equation}\label{eq:fp-operator}
  \Ladj u = -\sum_i \partial_{x_i}\!\big(f_i\,u\big)
            + \tfrac12\sum_{i} \partial^{2}_{x_i}\!\big(\Sigma_{ii}\,u\big).
\end{equation}
For our model the diffusion is diagonal and degenerate: only the bath modes
carry noise, with
\begin{equation}\label{eq:diffusion}
  \tfrac12\Sigma_{I_j I_j} = 4\gamma T_j I_j, \qquad
  \tfrac12\Sigma_{\varphi_j\varphi_j} = 4\gamma T_j/I_j,
  \qquad j\in\{1,n\},
\end{equation}
and zero otherwise. The factor $4$ (rather than $2$) in \eqref{eq:diffusion}
reflects that the measured angles are doubled, $\theta_j = 2(\varphi_j -
\varphi_{j-1})$, so the angular noise variance is multiplied by $2^2$.

\begin{remark}[Degeneracy / hypoellipticity]\label{rem:hypo}
Because noise enters only at $j\in\{1,n\}$, $\Sigma$ has rank $4$ in the
$2n$-dimensional state space; the interior modes are forced only through the
Hamiltonian coupling. Existence, uniqueness and ergodicity of the invariant
measure for the $n=3$ system are established in \cite{HLNS} via a
Feynman--Kac--Lyapunov argument together with H\"ormander's condition; the long
chain $n>3$ is left there as an open problem and is the subject of
\cref{sec:long-chain}.
\end{remark}

The stationary density solves $\Ladj u = 0$ with $\int u\dd x = 1$. The
long-time relaxation toward it is governed by the (backward) generator $\Lgen$,
the formal adjoint of $\Ladj$,
\begin{equation}\label{eq:generator}
  \Lgen g = \sum_i f_i\,\partial_{x_i} g
          + \tfrac12\sum_i \Sigma_{ii}\,\partial^2_{x_i} g,
  \qquad
  (\E_x[g(X_t)]) = (e^{t\Lgen} g)(x).
\end{equation}
Under ergodicity the spectrum of $\Lgen$ lies in $\{\mathrm{Re}\,\lambda\le 0\}$
with a simple eigenvalue $\lambda_0=0$ (constant eigenfunction). Ordering the
remaining eigenvalues by real part, the \emph{spectral gap}
$\lambda_R := \max\{\mathrm{Re}\,\lambda : \lambda\neq 0\}<0$ sets the slowest
relaxation rate, and for an observable $g$ orthogonal to constants
\begin{equation}\label{eq:decay}
  \E_x[g(X_t)] - \E_{\mathrm{eq}}[g]
  \;\sim\; Q(x)\,e^{\lambda_R t}\quad (t\to\infty),
\end{equation}
where $Q$ is the corresponding (sub)dominant eigenfunction, $\Lgen Q =
\lambda_R Q$. A purely real $\lambda_R$ indicates an overdamped relaxation; a
nonzero imaginary part $\lambda_I$ indicates damped oscillation with frequency
$\lambda_I$.

\subsection{Numerical methods}\label{sec:methods}

\subsubsection{A data-driven neural-network Fokker--Planck solver}
Classical grid-based PDE solvers are impractical for \eqref{eq:fp-operator} in
dimension $2n$: the domain is unbounded with no natural boundary condition, and
the cost grows exponentially with dimension. We follow the mesh-free,
data-driven solver of Zhai--Dobson--Li \cite{ZhaiDobsonLi} (building on
\cite{Li2019,DobsonLiZhai}), which represents the stationary density by a neural
network $\tilde u(x;\vartheta)$ trained against the unconstrained loss
\begin{equation}\label{eq:nn-loss}
  \mathcal{J}(\vartheta)
  = \underbrace{\frac{1}{N_X}\sum_{i}\big(\Ladj \tilde u(x_i;\vartheta)\big)^2}_{\text{PDE residual } L_1}
  + \underbrace{\frac{1}{N_Y}\sum_{j}\big(\tilde u(y_j;\vartheta)-v(y_j)\big)^2}_{\text{MC reference } L_2},
\end{equation}
where $\Ladj\tilde u$ is evaluated by automatic differentiation and $v$ is a
low-accuracy Monte Carlo (MC) estimate of the density at reference points
$\{y_j\}$. The operator term $L_1$ regularizes the network toward solutions of
$\Ladj u=0$ and sharply reduces the amount and accuracy of MC data required;
the method tolerates spatially uncorrelated noise in $v$ \cite{ZhaiDobsonLi}.

We adapt the solver to \eqref{eq:sde-bulk} in two ways. First, the angles are
periodic, so the network input is the \emph{periodic encoding}
$(\bI,\cos\theta_1,\sin\theta_1,\dots)$ rather than the raw angles. Second, the
diffusion \eqref{eq:diffusion} is degenerate (\cref{rem:hypo}); we therefore
include collocation points in the zero-density region and constrain them by the
operator $L_1$ alone (no fabricated reference data there). The reference density
$v$ is obtained by a high-throughput MC simulation of \eqref{eq:sde-bulk}
together with a kernel density estimate; see \cref{sec:fp-numerics}.

\subsubsection{Eigenfunctions of the generator}
To probe \eqref{eq:decay} we estimate the dominant nontrivial eigenpair of
$\Lgen$ directly from trajectories. For a localized observable $g$ (we use
$g=\cos\theta_1$) we evolve a large ensemble of trajectories from each of many
initial points $x$ and record the conditional expectation $\E_x[g(X_t)]$ as a
function of $t$. Fitting the tail \eqref{eq:decay} yields the spectral gap
$\lambda_R$ (and, via the presence or absence of oscillation, $\lambda_I$),
while the fitted amplitude gives a pointwise sample $Q(x)$ of the eigenfunction.
A second network is then trained to represent $Q$ on the whole domain, using a
data loss against the sampled $Q(x)$ together with a PDE loss enforcing the
eigenvalue equation $\Lgen Q=\lambda_R Q$. Results are reported in
\cref{sec:eigenfunction}.

\section{Statistical physics properties of the long cascade chain}\label{sec:long-chain}

The rigorous analysis of \cite{HLNS} treats the three-mode chain and lists the
longer chains $n>3$ as open. This section studies those long chains
numerically: the structure of the nonequilibrium steady state
(\cref{sec:ness-long}), heat transport and the validity of Fourier's law
(\cref{sec:thermal-conductivity}), and entropy production
(\cref{sec:entropy}).

\paragraph{Simulation.}
All long-chain data are produced by a single Euler--Maruyama integrator of
\eqref{eq:sde-bulk} that processes $W=16$ independent trajectories in lockstep
(one SIMD register per arithmetic operation), with the phases wrapped to
$[-\pi,\pi)$ at every step. We use $\gamma=0.1$, time step $\Delta t=10^{-3}$,
a burn-in of length $T_{\mathrm{burn}}$, and sample every $100$ steps to keep
successive samples approximately decorrelated. The \emph{closed chain} (denoted
case~$0$) couples both end modes to baths through \eqref{eq:bath}; the
\emph{open chain} (case~$4$) leaves the cold end free. Three-dimensional
histograms of $(I_j,I_{j+1},\theta_j)$, their one-dimensional marginals, and the
mean action profile are accumulated on the fly, so that runs of order $10^9$
samples are feasible without storing trajectories.

\subsection{Nonequilibrium steady state}\label{sec:ness-long}

\paragraph{Stabilization.}
We regard a run as stabilized when the sampled expectations (in particular the
action profile $\langle I_j\rangle$ and the bond-angle statistics) cease to
change with increasing sample size. Because the cascade transports energy
slowly, the relaxation time is long and the burn-in must be chosen with care:
a short run can report a spuriously uniform profile that has not yet developed
the steady gradient. We return to this point below.

\paragraph{Energy profile.}
\Cref{fig:profile} shows the mean action profile $\langle I_j\rangle$ for the
closed chain with $T_1=10$, $T_n=2$ at $n=25$. A long run
($T_{\mathrm{final}}=2\times10^{5}$, $\sim5\times10^{8}$ samples) exhibits a
clear monotone profile, decreasing smoothly from $\langle I_1\rangle\approx0.74$
at the hot end to $\langle I_{24}\rangle\approx0.39$ in the bulk, with a sharp
drop to $\langle I_{25}\rangle\approx0.17$ at the cold contact. A much shorter
run ($T_{\mathrm{final}}=3\times10^{4}$) instead reports an almost flat profile
$\langle I_j\rangle\approx0.64$; we attribute the discrepancy to insufficient
equilibration in the short run rather than to a genuine plateau, and the
profiles reported here use the long, converged runs.
% TODO: confirm with a burn-in convergence study; add profiles at several N and
%       the open-chain (case 4) profile on the same axes.

\begin{figure}[ht]
  \centering
  % \includegraphics[width=0.9\textwidth]{thermal_profile.png} % TODO: insert PNG
  \fbox{\parbox{0.7\textwidth}{\centering\textit{[figure placeholder]}\\
  Mean action profile $\langle I_j\rangle$ vs.\ mode index $j$, closed chain,
  $T_1=10,T_n=2$, $n=25$ (linear and log scales); short- vs.\ long-run
  comparison.}}
  \caption{Steady-state action profile of the long closed chain.}
  \label{fig:profile}
\end{figure}

\paragraph{Terminal energy and the cascade geometry.}
For the open chain one expects the action to settle, far from the driven end, to
a length-dependent \emph{terminal energy}; tracking this terminal value as $n$
grows probes how deep the cascade penetrates. The interpretation connects the
chain index to genuine Fourier modes through the frequency embedding of
\cite{CKSTT}: each generation of the resonant ``nuclear family'' construction
places a group of modes on $\mathbb{Z}^2$ by exact rational (Pythagorean)
rotations, so that the abstract chain position $j$ corresponds to a concrete set
of frequencies $\Lambda_j$. \Cref{fig:embedding} plots the locations of the
first several groups, computed with exact rational arithmetic.
% TODO(data): terminal-energy-vs-N curve for the open chain (case 4).
% Embedding figure available from the rational-arithmetic construction.

\begin{figure}[ht]
  \centering
  % \includegraphics[width=0.7\textwidth]{combined_plot_N5_R100.png} % TODO: insert PNG
  \fbox{\parbox{0.7\textwidth}{\centering\textit{[figure placeholder]}\\
  Locations on $\mathbb{Z}^2$ of the first generations of the CKSTT frequency
  embedding; (inset/TODO) terminal energy vs.\ $n$ for the open chain.}}
  \caption{Cascade geometry: frequency embedding and terminal energy.}
  \label{fig:embedding}
\end{figure}

\paragraph{Local thermodynamic equilibrium.}\label{par:lte}
A central question is whether, despite the global nonequilibrium drive, adjacent
modes in the bulk are described \emph{locally} by a Gibbs-type measure. We test
this at an interior pair $(j,j+1)$ of the $n=25$ chain ($T_1=10,T_n=2$) by
accumulating the joint density $\rho(I_j,I_{j+1},\theta_j)$ over $\sim5\times
10^{8}$ samples and fitting $\log\rho$ to three candidate forms:
\begin{align}
  \text{(A)}\quad &\log\rho \sim a + b\big[-\tfrac14(I_a^2+I_b^2)+I_a I_b\cos\theta\big],
  \label{eq:lteA}\\
  \text{(B)}\quad &\log\rho \sim a + b\big[\tfrac14(I_a^2+I_b^2)+I_a I_b(1+\cos\theta)\big],
  \label{eq:lteB}\\
  \text{(C)}\quad &\log\rho \sim a + c_1 I_a^2 + c_2 I_b^2 + c_3 I_a I_b + c_4 I_a I_b\cos\theta,
  \label{eq:lteC}
\end{align}
where $(I_a,I_b)=(I_j,I_{j+1})$. Form~(B) is a local Gibbs measure at an
effective inverse temperature $-b$; form~(C) is the general second-order ansatz
that lets the data choose. The naive guess (A) fails badly
($R^2\approx0.27$), whereas the local-Gibbs form (B) captures most of the
structure ($R^2\approx0.87$) and the general form (C) does slightly better
($R^2\approx0.91$); see \cref{fig:lte-fit,fig:lte-slices}. Thus the bulk does
satisfy a local thermodynamic equilibrium, but of local-Gibbs type
\eqref{eq:lteB}, not the form originally conjectured \eqref{eq:lteA}.

\begin{figure}[ht]
  \centering
  % \includegraphics[width=0.95\textwidth]{lte_fit_j12.png} % TODO: insert PNG
  \fbox{\parbox{0.8\textwidth}{\centering\textit{[figure placeholder]}\\
  Predicted vs.\ observed $\log\rho$ for forms (A), (B), (C) at the central site,
  $R^2=0.27/0.87/0.91$.}}
  \caption{Local-equilibrium fits of the joint pair density.}
  \label{fig:lte-fit}
\end{figure}

A scan across all interior sites confirms the picture and refines it. The fitted
coefficients of \eqref{eq:lteC} satisfy $c_1\approx c_2$ (symmetry between the
two neighboring actions) with $R^2\approx0.90$ at every interior site. The
angle coupling, however, is \emph{suppressed} relative to a strict local Gibbs
measure: a local Gibbs form predicts $c_3=c_4$, whereas the data give
$c_3/c_4\approx1.9$ in the bulk (rising toward the contacts). Correspondingly
the one-dimensional bond-angle marginal is bimodal, with maxima near
$\theta=\pm\pi$ and a minimum at $\theta=0$ (\cref{fig:lte-marg}).
% TODO: regenerate the multi-site scan (c1,c2,c3,c4 and implied local T) on the
%       long, converged histograms (the SIMD accumulator now outputs marginals);
%       reconcile the implied local-temperature profile with Fig. profile.

\begin{figure}[ht]
  \centering
  % \includegraphics[width=0.95\textwidth]{lte_marginals_j12.png} % TODO: insert PNG
  \fbox{\parbox{0.8\textwidth}{\centering\textit{[figure placeholder]}\\
  One-dimensional marginals of $I_j$, $I_{j+1}$ (near-exponential) and
  $\theta_j$ (bimodal, peaks near $\pm\pi$).}}
  \caption{Marginals at an interior site.}
  \label{fig:lte-marg}
\end{figure}

\begin{figure}[ht]
  \centering
  % \includegraphics[width=0.95\textwidth]{theta_slices_empirical_j12.png} % TODO
  \fbox{\parbox{0.8\textwidth}{\centering\textit{[figure placeholder]}\\
  Conditional density $\rho(I_a,I_b\mid\theta)$ at $\theta=-\pi,-\pi/2,0,\pi/2$
  with the contours of forms (A) and (B) overlaid.}}
  \caption{Angle-resolved test of local equilibrium.}
  \label{fig:lte-slices}
\end{figure}

\subsection{Thermal conductivity}\label{sec:thermal-conductivity}

In the closed-chain NESS the mean action current $\E[J]$ of
\eqref{eq:continuity} is uniform along the chain. We measure how it scales with
the chain length $n$. A least-squares fit over $n\in[5,100]$ gives
\begin{equation}\label{eq:flux-scaling}
  \E[J] \;\sim\; n^{-1.75}\qquad (R^2\approx0.92),
\end{equation}
so that $\E[J]\cdot n\to0$: the current decays faster than $1/n$. By
\cref{de:fourier} the effective conductivity behaves as $\kappa_n\sim n^{-0.75}
\to 0$, i.e.\ the chain is a sub-diffusive thermal \emph{insulator}---the
opposite of the divergent conductivity familiar from FPU-type lattices
\cite{LepriLiviPoliti}. \Cref{fig:flux-scaling} shows the fit.
% CAVEAT: the exponent was extracted from data generated with an earlier
% integrator that carried a boundary phi-drift sign error; re-confirm on the
% canonical SIMD code before quoting the exponent as final.

\begin{figure}[ht]
  \centering
  % \includegraphics[width=0.7\textwidth]{scaling.png} % TODO: insert PNG
  \fbox{\parbox{0.7\textwidth}{\centering\textit{[figure placeholder]}\\
  Mean current $\E[J]$ vs.\ chain length $n$ (log--log), with the fit
  $\E[J]\sim n^{-1.75}$.}}
  \caption{Anomalous scaling of the heat current.}
  \label{fig:flux-scaling}
\end{figure}

As a consistency check, the open chain (one free end) carries zero net steady
current, $\E[J]=0$ within statistical error. We also record the dependence of
$\E[J]$ on the imposed temperature difference $\Delta T=T_1-T_n$ at fixed $n$;
the current is non-monotone in $\Delta T$, rising and then saturating, with a
maximum near $\Delta T\approx 8$.
% TODO(data): flux-vs-Delta-T curve (figure); current is non-monotone, peak ~8-9.

Beyond the mean, the \emph{distribution} of the instantaneous current is of
interest. We expect an exponential large-deviation tail,
\begin{equation}\label{eq:flux-tail}
  \mathrm{Prob}[\,J>C\,]\;\approx\;e^{-rC}\qquad (C\ \text{large}),
\end{equation}
for some rate $r>0$. Testing \eqref{eq:flux-tail} requires recording the
instantaneous current $J(t)$ along the trajectory (rather than only its mean),
which is a small extension of the accumulator.
% TODO(data): record instantaneous J(t); histogram and fit the tail rate r.

\subsection{Entropy production and its fluctuation}\label{sec:entropy}

The two-bath NESS produces entropy at the rate \eqref{eq:entropy-prod}. Along a
trajectory the heat delivered by each bath is the work done by its
dissipative--stochastic force on the corresponding end mode; accumulating the
mean rates $\E[\dot Q_1],\E[\dot Q_n]$ yields the entropy production rate
$\dot S=\E[\dot Q](1/T_n-1/T_1)\ge0$, a direct numerical check of the second law
for this model. Recording instead the entropy produced over sliding windows of
length $\tau$ gives the empirical distribution of $\Sigma_\tau$, whose symmetry
under $\Sigma_\tau\to-\Sigma_\tau$ tests the Gallavotti--Cohen fluctuation
theorem \eqref{eq:ft}.
% TODO(data): no simulation material yet for this subsection. Implement the
% bath-heat-flux accumulation (Q_1, Q_n) and the windowed Sigma_tau statistics;
% report (i) the second-law check Sdot >= 0 and (ii) the FT slope test.


\section{Detailed property of the short chain}\label{sec:short-chain}

We now turn to the smallest nontrivial cascade, the three-mode system $n=3$,
for which existence, uniqueness and ergodicity of the NESS are rigorously known
\cite{HLNS}. The Hamiltonian and bath of \cref{sec:long-chain} depend on the
phases only through $\theta_1=2(\varphi_1-\varphi_2)$ and
$\theta_3=2(\varphi_3-\varphi_2)$, so the dynamics reduces from the six
variables $(\bI,\bphi)$ to the five-dimensional state
$x=(I_1,I_2,I_3,\theta_1,\theta_3)$, in which the middle mode $I_2$ carries no
bath. This low dimension lets us solve the stationary Fokker--Planck equation
directly and interrogate the invariant measure (\cref{sec:fp-numerics}), the
mechanism that stabilizes the NESS (\cref{sec:stabilization}), and the spectrum
of the generator (\cref{sec:eigenfunction}).

\subsection{Numerical solution of the nonequilibrium steady state}\label{sec:fp-numerics}

In the reduced coordinates the stationary density solves $\Ladj u=0$, with
$\Ladj$ as in \eqref{eq:fp-operator}: the drift is the $n=3$ specialization of
\eqref{eq:sde-bulk}, and the (degenerate) diffusion \eqref{eq:diffusion} acts
only on $(I_1,\theta_1)$ and $(I_3,\theta_3)$ with coefficients
$4\gamma T_1 I_1,\ 4\gamma T_1/I_1$ and $4\gamma T_3 I_3,\ 4\gamma T_3/I_3$. We
solve it with the neural-network method of \cref{sec:methods}. The reference
density $v$ is produced by a long SIMD Monte Carlo integration of
\eqref{eq:sde-bulk} followed by a kernel density estimate on a grid of boxes in
$x$; boxes in which the Monte Carlo records no sample are retained as
zero-density collocation points and are constrained by the operator term $L_1$
alone. The network outputs the log-density, $u(x)$ with $\rho=\exp(\sigma_0
u+\mu_0)$ for fixed normalization constants $(\mu_0,\sigma_0)$, takes the
periodic encoding $(I_1,I_2,I_3,\cos\theta_1,\sin\theta_1,\cos\theta_3,
\sin\theta_3)$ as input, and is trained in two phases: a data-only pre-training
against $\{v(y_j)\}$, followed by joint training in which the weight of the PDE
residual $L_1$ is ramped on gradually.

\paragraph{Validation against the equilibrium Gibbs measure.}
When $T_1=T_3=T$ the exact invariant density is the Gibbs measure
\eqref{eq:gibbs}, which provides a stringent test. Training at $T=5$ and
comparing the network density to \eqref{eq:gibbs} on the slices $I_1=I_2=I_3=I$
(\cref{fig:eq-validation}), the mean relative error is $1.5\%$, $2.6\%$ and
$6.4\%$ at $I=0.5,1,2$ respectively, with the pointwise $\log$-ratio
$\log(\rho_{\mathrm{NN}}/\rho_{\mathrm{Gibbs}})$ confined to a narrow band. The
solver thus recovers the known equilibrium measure across the range of actions
that carry appreciable probability.

\begin{figure}[ht]
  \centering
  % \includegraphics[width=0.9\textwidth]{eq_validation.png} % TODO: insert PNG
  \fbox{\parbox{0.7\textwidth}{\centering\textit{[figure placeholder]}\\
  Network density vs.\ analytic Gibbs measure at $T_1=T_3=5$, slices $I=0.5,1,2$;
  relative error and $\log$-ratio.}}
  \caption{Equilibrium validation of the neural-network Fokker--Planck solver.}
  \label{fig:eq-validation}
\end{figure}

\paragraph{Nonequilibrium steady state.}
With distinct bath temperatures $T_1=2$, $T_3=8$ the invariant measure is no
longer of Gibbs form. \Cref{fig:neq-density} shows the network density on the
angular plane $(\theta_1,\theta_3)$ at several values of the common action: the
distribution concentrates around a phase-locked configuration and is visibly
distorted relative to the equilibrium measure. The departure from equilibrium
is quantified in \cref{sec:stabilization}.

\begin{figure}[ht]
  \centering
  % \includegraphics[width=0.9\textwidth]{neq_density.png} % TODO: insert PNG
  \fbox{\parbox{0.7\textwidth}{\centering\textit{[figure placeholder]}\\
  Network density $\rho(\theta_1,\theta_3)$ at $T_1=2,T_3=8$ for several $I$,
  with the equilibrium ($T_1=T_3=5$) density for comparison.}}
  \caption{Nonequilibrium steady-state angular density of the three-mode chain.}
  \label{fig:neq-density}
\end{figure}

\subsection{Mechanism of stabilization}\label{sec:stabilization}

The NESS is stabilized by a phase-locking mechanism that operates differently in
the high- and low-energy regimes \cite{HLNS}. At high energy the bond angles are
driven to a stable equilibrium near $2\pi/3$, escaping the unstable
configuration at $4\pi/3$; at low energy the angle $\theta_3$ remains near
$4\pi/3$, which prevents the middle action $I_2$ from collapsing to zero and
keeps the cascade alive. The reduced four-dimensional subsystem
$(I_1,I_3,\theta_1,\theta_3)$, obtained in the low-energy limit and equipped
with the $\gamma(T-I^3)$ thermostat of \cite{HLNS}, isolates this mechanism; we
solve its stationary Fokker--Planck equation by the same network method and
recover the predicted concentration of $\theta_3$.

Three quantitative signatures of the nonequilibrium drive are read off from the
$T_1=2,T_3=8$ steady state.

\emph{(i) Symmetry breaking.} At equilibrium the density is symmetric under
$\theta_1\leftrightarrow\theta_3$. Out of equilibrium this symmetry is broken:
the relative asymmetry $|\rho(\theta_1,\theta_3)-\rho(\theta_3,\theta_1)|/
(\rho(\theta_1,\theta_3)+\rho(\theta_3,\theta_1))$ reaches a maximum of about
$67\%$ and averages about $28\%$ over the angular plane (\cref{fig:symmetry}).

\emph{(ii) Angular width.} The two angles relax with widths set by their bath
temperatures: the ratio of circular standard deviations satisfies
$\sigma(\theta_3)/\sigma(\theta_1)\approx\sqrt{T_3/T_1}=2$ across the range of
actions, consistent with each angle behaving locally like an overdamped
oscillator thermalized to its own bath.

\emph{(iii) Phase-locking angle.} The most probable bond angle is shifted from
its equilibrium value to
\begin{equation}\label{eq:phase-lock}
  \theta_0 \;=\; \arcsin\!\frac{\sqrt{4\gamma^2+3}-\gamma}{2(\gamma^2+1)},
\end{equation}
which the network density reproduces at the measured value of $\gamma$.

\begin{remark}\label{rem:eq47}
The derivation of \eqref{eq:phase-lock} (the fixed point of the angular drift
balanced against dissipation) is to be supplied; the formula is recorded here so
the numerical comparison can be stated.
\end{remark}

Finally, current conservation through the unforced middle mode gives a balance
condition: in the NESS the two action currents into $I_2$ cancel in the mean,
$\E[\,I_2 I_1\sin\theta_1\,]+\E[\,I_2 I_3\sin\theta_3\,]\approx 0$, which the
steady-state density satisfies and which we use as an internal consistency check
on the solver.

\begin{figure}[ht]
  \centering
  % \includegraphics[width=0.9\textwidth]{symmetry_breaking.png} % TODO: insert PNG
  \fbox{\parbox{0.7\textwidth}{\centering\textit{[figure placeholder]}\\
  $\rho(\theta_1,\theta_3)$, its transpose, and the relative asymmetry at
  $T_1=2,T_3=8$.}}
  \caption{$\theta_1\leftrightarrow\theta_3$ symmetry breaking in the
  nonequilibrium steady state.}
  \label{fig:symmetry}
\end{figure}

\subsection{Eigenfunction of the generator}\label{sec:eigenfunction}

To probe the relaxation spectrum we apply the trajectory method of
\cref{sec:methods} with the localized observable $g=\cos\theta_1$. From each of
a large set of initial points we evolve an ensemble and record
$\E_x[\cos\theta_1(X_t)]$; fitting the decay \eqref{eq:decay} returns, at each
initial point, an estimate of the relaxation rate and amplitude.

At $\gamma=0.1$ the relaxation is \emph{overdamped}: the distribution of fitted
imaginary parts $\lambda_I$ is sharply peaked at zero, so the dominant part of
the spectrum is real (\cref{fig:eigenvalues}). The spectral gap, estimated as
the median of the pure-exponential fits, is
\begin{equation}\label{eq:lambdaR}
  \lambda_R \;\approx\; -0.93 .
\end{equation}
(An oscillatory five-parameter fit applied indiscriminately to all trajectories
returns a more negative mean, $\approx-1.90$; this mixes in faster modes and the
early transient and is reported only as a diagnostic confirming that most
trajectories do not oscillate.)

The fitted amplitudes provide pointwise samples $Q(x)$ of the dominant
nontrivial eigenfunction. A second network is trained to represent $Q$ subject
to the eigenvalue equation $\Lgen Q=\lambda_R Q$ with $\lambda_R$ from
\eqref{eq:lambdaR}, using the generator $\Lgen$ of \eqref{eq:generator} (without
the divergence terms that distinguish $\Ladj$). The reconstruction is only
moderately accurate: the data root-mean-square error is $0.38$ (relative error
$0.71$) and the median normalized PDE residual $|\Lgen Q-\lambda_R Q|/|Q|$ is
$0.35$. We report these candidly: the dominant nontrivial eigenfunction is
sign-changing and is reconstructed from intrinsically noisy backward Monte
Carlo data, and is the hardest object computed in this work; the agreement is
nonetheless sufficient to exhibit the qualitative structure of $Q$
(\cref{fig:eigenfunction}).

\begin{figure}[ht]
  \centering
  % \includegraphics[width=0.9\textwidth]{eigenvalue_scatter.png} % TODO: insert PNG
  \fbox{\parbox{0.7\textwidth}{\centering\textit{[figure placeholder]}\\
  Distributions of fitted $\lambda_R$ and $\lambda_I$ over initial conditions;
  the spike of $\lambda_I$ at $0$ indicates a real spectrum.}}
  \caption{Generator spectrum estimated from backward Monte Carlo at
  $\gamma=0.1$.}
  \label{fig:eigenvalues}
\end{figure}

\begin{figure}[ht]
  \centering
  % \includegraphics[width=0.9\textwidth]{Q1_slices.png} % TODO: insert PNG
  \fbox{\parbox{0.7\textwidth}{\centering\textit{[figure placeholder]}\\
  Network eigenfunction $Q(\theta_1,\theta_3)$ at $I=1,2,4$.}}
  \caption{Dominant nontrivial eigenfunction of the generator.}
  \label{fig:eigenfunction}
\end{figure}


\begin{thebibliography}{9}

\bibitem{CKSTT}
J.~Colliander, M.~Keel, G.~Staffilani, H.~Takaoka, T.~Tao,
\emph{Transfer of energy to high frequencies in the cubic defocusing nonlinear
Schr\"odinger equation}, Invent. Math. \textbf{181} (2010), 39--113.

\bibitem{HLNS}
Z.~Hani, Y.~Li, A.~R.~Nahmod, G.~Staffilani,
\emph{Nonequilibrium steady states for a three-mode energy cascade model},
arXiv:2505.16018.

\bibitem{ZhaiDobsonLi}
J.~Zhai, M.~Dobson, Y.~Li,
\emph{A deep learning method for solving Fokker--Planck equations},
Proc. Mach. Learn. Res. (MSML) \textbf{145} (2022), 568--597.

\bibitem{Li2019}
Y.~Li,
\emph{A data-driven method for the steady state of randomly perturbed dynamics},
Commun. Math. Sci. \textbf{17} (2019), no.~4, 1045--1059.

\bibitem{DobsonLiZhai}
M.~Dobson, Y.~Li, J.~Zhai,
\emph{An efficient data-driven solver for Fokker--Planck equations: algorithm
and analysis}, arXiv:1906.02600 (2019).

\bibitem{GallavottiCohen}
G.~Gallavotti, E.~G.~D.~Cohen,
\emph{Dynamical ensembles in nonequilibrium statistical mechanics},
Phys. Rev. Lett. \textbf{74} (1995), 2694--2697.

\bibitem{LepriLiviPoliti}
S.~Lepri, R.~Livi, A.~Politi,
\emph{Thermal conduction in classical low-dimensional lattices},
Phys. Rep. \textbf{377} (2003), 1--80.

\bibitem{Nazarenko}
S.~Nazarenko, \emph{Wave Turbulence}, Lecture Notes in Physics \textbf{825}, Springer, 2011.

\end{thebibliography}


\end{document}
